\chapter{Linear and Affine Functions}\label{ch:g9:linfunc}

``Three kilograms cost three times as much as one'': that is
proportionality, the simplest way two quantities can be related. Linear
functions are proportionality in the language of functions; affine
functions add a fixed starting amount. Their graphs are straight lines,
and reading those lines is the skill this chapter builds.

\section{Proportionality and linear functions}

\begin{definition}[Linear function]\label{def:g9:linfunc:linear}
A \emph{linear function}\index{linear function} multiplies every number
by a fixed number $a$:
\[
f(x) = a x .
\]
The number $a$ is the \emph{coefficient} of $f$. Two quantities are
\emph{proportional}\index{proportionality} exactly when one is a linear
function of the other.
\end{definition}

\begin{example}\label{ex:g9:linfunc:linear}
If apples cost $3$ per kilogram, the price of $x$ kilograms is
$f(x) = 3x$: buying twice as much costs twice as much, and the price of
$2.5$ kg is $f(2.5) = 7.5$.

Conversely, if a linear function satisfies $f(4) = 10$, then
$a = \frac{10}{4} = 2.5$ and $f(x) = 2.5x$: \emph{one value determines a
linear function completely.}
\end{example}

\begin{proposition}[Graph of a linear function]\label{prop:g9:linfunc:lineargraph}
The graph of $f(x) = ax$ is a straight line \emph{through the origin}.
Conversely, every non-vertical line through the origin is the graph of a
linear function.
\end{proposition}
\admitted

\begin{omfigure}
\begin{tikzpicture}
\begin{axis}[omaxis,
  width=7.4cm, height=5.6cm,
  xmin=-3.4, xmax=3.4, ymin=-3.2, ymax=3.6,
  xtick={-2,-1,1,2}, ytick={-2,-1,1,2}]
  \addplot[omDef, thick, domain=-2.2:2.2] {1.5*x};
  \addplot[omThm, thick, domain=-3.1:3.1] {-0.5*x};
  \node[omDef, font=\small] at (axis cs:2.15,2.6) {$y = 1.5x$};
  \node[omThm, font=\small] at (axis cs:-2.5,1.7) {$y = -0.5x$};
\end{axis}
\end{tikzpicture}

{\small Two linear functions: both graphs pass through the origin; the
coefficient $a$ is the height reached at $x = 1$.}
\end{omfigure}

\section{Affine functions}

\begin{definition}[Affine function]\label{def:g9:linfunc:affine}
An \emph{affine function}\index{affine function} has the form
\[
f(x) = a x + b ,
\]
where $a$ and $b$ are fixed numbers. Linear functions are the special
case $b = 0$; constant functions the case $a = 0$.
\end{definition}

\begin{example}
A gym charges a registration fee of $20$, then $5$ per visit: the total
cost of $x$ visits is $f(x) = 5x + 20$. The cost is \emph{not}
proportional to the number of visits (10 visits do not cost twice as much
as 5), but each extra visit adds the same amount, $5$.
\end{example}

\begin{proposition}[Graph of an affine function]\label{prop:g9:linfunc:affinegraph}
The graph of $f(x) = ax + b$ is a straight line:
\begin{itemize}
  \item $b = f(0)$ is the \emph{$y$-intercept}: the line crosses the
        vertical axis at $(0, b)$;
  \item $a$ is the \emph{slope}\index{slope}: moving one unit to the
        right moves $a$ units up (down if $a < 0$); the function is
        increasing when $a > 0$, decreasing when $a < 0$;
  \item for any two inputs $u \neq v$,
        $a = \dfrac{f(v) - f(u)}{v - u}$.
\end{itemize}
\end{proposition}

\begin{proof}[Proof of the last point]
$f(v) - f(u) = (av + b) - (au + b) = a(v - u)$; divide by $v - u$. The
rest is admitted at this level (see \cref{ch:g10:lines} for a full
treatment of lines and their equations).
\end{proof}

\begin{omfigure}
\begin{tikzpicture}
\begin{axis}[omaxis,
  width=7.8cm, height=5.8cm,
  xmin=-1.9, xmax=4.4, ymin=-1.4, ymax=5.4,
  xtick={1,2,3}, ytick={1,2,3}]
  \addplot[omDef, thick, domain=-1.6:4.1] {x + 1};
  \fill (axis cs:0,1) circle (1.5pt);
  \node[left, font=\small] at (axis cs:0,1.25) {$b$};
  \draw[dashed, omThm] (axis cs:1,2) -- (axis cs:2,2) -- (axis cs:2,3);
  \node[font=\footnotesize, below] at (axis cs:1.5,2) {$+1$};
  \node[font=\footnotesize, right] at (axis cs:2,2.5) {$+a$};
\end{axis}
\end{tikzpicture}

{\small Reading $y = ax + b$ on its graph: the line crosses the $y$-axis
at height $b$, and climbs $a$ for each unit step to the right (here
$a = 1$, $b = 1$).}
\end{omfigure}

\begin{method}[Finding an affine function from two values]\label{met:g9:linfunc:findaffine}
Given $f(u)$ and $f(v)$ with $u \neq v$:
\begin{enumerate}
  \item compute the slope $a = \dfrac{f(v) - f(u)}{v - u}$;
  \item substitute one of the two known values into $f(x) = ax + b$ to
        find $b$;
  \item check with the other value.
\end{enumerate}
\end{method}

\begin{example}\label{ex:g9:linfunc:findaffine}
Find the affine function with $f(2) = 7$ and $f(5) = 16$.
\[
a = \frac{16 - 7}{5 - 2} = \frac93 = 3 ,
\]
then $f(2) = 3 \times 2 + b = 7$ gives $b = 1$: $f(x) = 3x + 1$. Check:
$f(5) = 15 + 1 = 16$.
\end{example}

\section{Percentages}

\begin{proposition}[Percentage change]\label{prop:g9:linfunc:percent}
Increasing a quantity by $t\,\%$ multiplies it by
$1 + \dfrac{t}{100}$; decreasing it by $t\,\%$ multiplies it by
$1 - \dfrac{t}{100}$. Percentage changes are linear functions.
\end{proposition}

\begin{proof}
Increasing $x$ by $t\,\%$ means adding $\frac{t}{100}\,x$:
\[
x + \frac{t}{100}\,x = \left(1 + \frac{t}{100}\right) x .
\]
Same computation with a minus sign for a decrease.
\end{proof}

\begin{example}\label{ex:g9:linfunc:percent}
A price of $80$ increases by $15\%$: new price
$80 \times 1.15 = 92$. A price of $60$ decreases by $30\%$:
$60 \times 0.7 = 42$.

\emph{Chaining changes multiplies the factors.} A $20\%$ increase
followed by a $20\%$ decrease gives
$1.2 \times 0.8 = 0.96$: a $4\%$ \emph{decrease} overall, not a return to
the start!
\end{example}

\section{Exercises}

\begin{exercise}[$\star$]\label{exo:g9:linfunc:1}
Among these functions, which are linear? affine? neither?
\[
f(x) = 4x, \qquad
g(x) = 3x - 5, \qquad
h(x) = x^2, \qquad
k(x) = -\tfrac x2, \qquad
m(x) = 7 .
\]
\end{exercise}

\begin{exercise}[$\star$]\label{exo:g9:linfunc:2}
Let $f(x) = 2x - 3$. Compute $f(0)$, $f(4)$, $f(-1)$, and find the number
$x$ such that $f(x) = 9$.
\end{exercise}

\begin{exercise}[$\star$]\label{exo:g9:linfunc:3}
A linear function satisfies $f(6) = 15$. Find its coefficient, then
compute $f(10)$ and $f(-2)$.
\end{exercise}

\begin{exercise}[$\star$]\label{exo:g9:linfunc:4}
Draw on the same coordinate system the graphs of $f(x) = 2x$,
$g(x) = 2x + 3$ and $h(x) = -x + 4$. What can be said about the graphs of
$f$ and $g$?
\end{exercise}

\begin{exercise}[$\star$]\label{exo:g9:linfunc:5}
Find the affine function such that $f(1) = 5$ and $f(3) = 11$; then the
one such that $g(0) = 4$ and $g(2) = 0$.
\end{exercise}

\begin{exercise}[$\star\star$]\label{exo:g9:linfunc:6}
Compute: the price after a $25\%$ increase on $120$; the price after a
$40\%$ discount on $65$; the original price if an article costs $69$
after a $25\%$ discount.
\end{exercise}

\begin{exercise}[$\star\star$]\label{exo:g9:linfunc:7}
A phone plan A costs $10$ per month plus $0.05$ per minute of calls; plan
B costs $25$ per month with unlimited calls.
\begin{enumerate}
  \item Express the monthly cost of plan A as a function of the number
        $x$ of minutes.
  \item From how many minutes per month is plan B cheaper?
\end{enumerate}
\end{exercise}

\begin{exercise}[$\star\star$]\label{exo:g9:linfunc:8}
A population of $2000$ bacteria increases by $10\%$ every hour.
\begin{enumerate}
  \item How many bacteria are there after one hour? After two hours?
  \item Explain why the answer after two hours is not $2400$.
\end{enumerate}
\end{exercise}

\begin{exercise}[$\star\star$]\label{exo:g9:linfunc:9}
The graph of an affine function passes through the points $(2, 3)$ and
$(6, 1)$.
\begin{enumerate}
  \item Compute its slope. Is the function increasing or decreasing?
  \item Give its expression, and compute the input for which the output
        is $0$.
\end{enumerate}
\end{exercise}

\begin{exercise}[$\star\star\star$]\label{exo:g9:linfunc:10}
A shop first raises a price by $t\,\%$, then lowers the new price by
$t\,\%$.
\begin{enumerate}
  \item Show that the final price equals the original multiplied by
        $1 - \dfrac{t^2}{10000}$.
  \item Deduce that the final price is always \emph{lower} than the
        original (for $0 < t \leq 100$), and find the overall percentage
        change for $t = 10$.
\end{enumerate}
\end{exercise}
